\chapter{Background}\label{ch:background}

In this chapter, we review the ideas that the rest of the thesis builds on. Sections~\ref{sec:bg-icl} to~\ref{sec:bg-linreg} introduce in-context learning, its interpretation as Bayesian inference, and in-context linear regression, a setting in which that interpretation can be computed exactly. Section~\ref{sec:bg-regret} defines the squared error by which we compare learners and gives the error of the Bayesian predictor. Sections~\ref{sec:bg-ess} and~\ref{sec:bg-mixture} present two tools from Bayesian statistics that we later apply to descriptions: the effective sample size of a prior and the robust mixture prior. Finally, Sections~\ref{sec:bg-theory} to~\ref{app:related} review the work most closely related to ours.

\section{In-Context Learning}\label{sec:bg-icl}
A language model takes a sequence of tokens, the prompt, and predicts what comes next. \citet{brown2020} showed that large language models can perform a task specified in the prompt, without any update to their weights. This ability is called in-context learning (ICL). \citet{brown2020} distinguished three kinds of prompts. A zero-shot prompt gives only a natural-language description of the task. One-shot and few-shot prompts add one or several \emph{demonstrations}, solved examples of the task, before the query that the model must answer. We follow their terminology. A demonstration is one input together with its label, and apart from the query we call everything else in the prompt a \emph{description}, whether it is written as a statement about the task or as a command. A prompt can contain both, and our question is how much the description adds to the demonstrations.

How models use the demonstrations is itself an active question. Pretrained models can perform well even when the labels of the demonstrations are replaced by random ones \citep{min2022}, which suggests that they use demonstrations partly to recognize a task they already know. Larger models, however, follow flipped in-context labels more readily \citep{wei2023}. \citet{lin2024} explain both behaviors, retrieving a known task and learning a new one, within a single Bayesian model. A description raises the same question from the other side. Does a model treat it as information about the task, to be weighed against the demonstrations, or does it use it in some other way?

\section{In-Context Learning as Bayesian Inference}\label{sec:bg-bayes}
Theoretical work studies ICL as Bayesian inference over latent tasks \citep{xie2022,panwar2024}, as a learning algorithm implemented in a forward pass \citep{akyurek2023,vonoswald2023}, and as meta-learning across tasks \citep{garg2022,muller2022}. Meta-trained predictors can approximate Bayesian inference under suitable training conditions \citep{genewein2025}, and the learning algorithms that Transformers implement in context approach the Bayes-optimal estimator as the networks grow \citep{akyurek2023}. Agreement can change with training-task diversity or new function classes \citep{raventos2023,panwar2024}, and language-model ICL can violate Bayesian consistency \citep{falck2024}. We adopt the Bayesian view, because it extends naturally from demonstrations to descriptions (Chapter~\ref{ch:setup}).

\citet{xie2022} proposed that ICL can be understood as implicit Bayesian inference. Suppose that pretraining text is generated from latent concepts. A model that predicts such text well must then infer the concept behind each document, and a prompt of demonstrations is evidence about the concept behind the prompt. In this view, a task is a latent parameter $w$, and a prompt contains $n$ demonstrations $D_n=\{(x_i,y_i)\}_{i=1}^n$ and a query $x$. The ideal learner predicts the query's label with the posterior predictive distribution
\begin{equation}\label{eq:bg-predictive}
    f(y \mid x, D_n) = \int f(y \mid x, w)\, f(w \mid D_n)\, dw .
\end{equation}

This view becomes testable when the prior over tasks is known. Consider a Transformer trained on prompts whose tasks are sampled from a known prior. Training on such prompts is called meta-training: the network is trained on many tasks and must infer from the context which task a prompt comes from. The posterior predictive minimizes the expected log loss over the prior (Section~\ref{sec:bg-regret}), and this expected log loss is exactly the training objective. A well-trained network therefore approximates the posterior predictive, whose mean is the Bayes-optimal prediction under squared error (Section~\ref{sec:bg-regret}) \citep{muller2022}. For regression with a Gaussian task prior, meta-trained Transformers behave like the Bayesian predictor, which is ridge regression \citep{akyurek2023, raventos2023, panwar2024}, provided that they are trained on enough distinct tasks \citep{raventos2023}. The Bayesian predictor is then a benchmark against which a network can be measured, and Chapter~\ref{ch:networks} uses it in this way.

\section{In-Context Linear Regression}\label{sec:bg-linreg}
\citet{garg2022} introduced in-context linear regression as a controlled test of ICL. Each prompt has its own task vector $w \sim \mathcal{N}(0, I_d)$, its inputs are drawn as $x_i \sim \mathcal{N}(0, I_d)$, and its labels are $y_i = w^\top x_i$. \citet{garg2022} found that a Transformer trained from scratch on such prompts predicts the query's label about as well as least squares. Later work added label noise, $y_i = w^\top x_i + \epsilon_i$ \citep{akyurek2023, raventos2023}. \citet{raventos2023} also varied the number of distinct tasks seen in training. With few tasks, the network behaves like the Bayesian predictor for a prior over those tasks alone. With enough tasks, it behaves like ridge regression, the Bayesian predictor for the Gaussian prior.

With Gaussian noise and a Gaussian prior, Bayes-optimal prediction has a closed form \citep[Section~14.8]{gelman2013}. We measure labels in units of their noise, so that $y_i \sim \mathcal{N}(w^\top x_i, 1)$. Let the prior be $w \sim \mathcal{N}(m, \Lambda^{-1})$, with center $m$ and precision $\Lambda$. Let $X_n \in \R^{n \times d}$ contain the demonstration inputs as rows, and write $W_n = X_n^\top X_n$. The posterior over $w$ is Gaussian, with covariance and mean
\begin{equation}\label{eq:bg-posterior}
    \Sigma = \big(\Lambda + W_n\big)^{-1}, \qquad \mu = \Sigma \big(\Lambda m + X_n^\top y_{1:n}\big),
\end{equation}
and the posterior predictive for the query is
\begin{equation}\label{eq:bg-gauss-predictive}
    y \mid x, D_n \sim \mathcal{N}\big(\mu^\top x,\; 1 + x^\top \Sigma x\big).
\end{equation}
The predictive variance has two parts. The first, $1$, is the label noise. The second, $x^\top \Sigma x$, is the remaining uncertainty about the task in the direction of the query. Under the baseline prior $\mathcal{N}(0, a_0 I_d)$, the posterior mean $\mu = (W_n + a_0^{-1} I_d)^{-1} X_n^\top y_{1:n}$ is the ridge regression estimate with penalty $a_0^{-1}$. Equation~\eqref{eq:bg-posterior} also shows that a prior acts like additional data. Each demonstration adds $x_i x_i^\top$ to the posterior precision, and the prior adds $\Lambda$. Chapter~\ref{ch:setup} uses this fact to give a description a prior of its own, and in Chapter~\ref{ch:specificity} a reliable description replaces the baseline precision $a_0^{-1}$ with $c$.

\section{Squared Error}\label{sec:bg-regret}
A predictor sees the prompt and gives a number $\hat y$ for the query's label. We score it by squared error. The label is $y = w^\top x + \epsilon$, where the noise $\epsilon$ has mean $0$ and variance $1$ and is independent of the prompt and the task. The expected squared error therefore splits into two parts:
\begin{equation}\label{eq:bg-sq}
    \E\big[(\hat y - y)^2\big] = \E\big[(\hat y - w^\top x)^2\big] + 1 .
\end{equation}
The second part is the noise variance, which is the same for every predictor. Following \citet{garg2022}, we therefore score a predictor by the first part alone, its \emph{prediction error}, or \emph{risk},
\[
    R = \E\big[(\hat y - w^\top x)^2\big],
\]
the expected squared distance from the noiseless label $w^\top x$. A predictor that knows the task predicts $w^\top x$ and has zero error.

Given the prompt, the task is uncertain. For any prediction $\hat y$ computed from the prompt,
\[
    \E\big[(\hat y - w^\top x)^2 \mid x, D_n\big] = \big(\hat y - \E[w^\top x \mid x, D_n]\big)^2 + \operatorname{Var}(w^\top x \mid x, D_n).
\]
The first term vanishes only when $\hat y$ is the posterior mean of $w^\top x$. When tasks are drawn from the prior, the posterior mean therefore has the smallest error of any predictor. We call it the Bayesian predictor and its error Bayes-optimal. It is also the mean of the posterior predictive~\eqref{eq:bg-predictive}. Its error is the expected posterior variance of $w^\top x$. In linear regression, the posterior is Gaussian with covariance $\Sigma$ (Equation~\eqref{eq:bg-posterior}), so the posterior variance of $w^\top x$ is $x^\top \Sigma x$, the part of the predictive variance~\eqref{eq:bg-gauss-predictive} that comes from uncertainty about the task. The query input is independent of the demonstrations, with $\E[x x^\top] = I_d$, so the error of the Bayesian predictor is
\begin{equation}\label{eq:bg-gauss-regret}
    R = \E\big[x^\top \Sigma x\big] = \E \operatorname{tr} \Sigma ,
\end{equation}
the expected trace of the posterior covariance. The error is zero when the posterior has concentrated on the task, and it grows with the uncertainty about the task that remains.

Two standard facts give this trace in the cases that later chapters need. Chapter~\ref{ch:setup} writes $R^{\mathrm{demo}}(n)$ for the error with $n$ demonstrations and no description. With fewer demonstrations than dimensions, $n < d$, the demonstrations say nothing about $d - n$ directions, so under the baseline prior $\mathcal{N}(0, a_0 I_d)$ we have $R^{\mathrm{demo}}(n) \ge (d - n)\,a_0$. Under a flat prior, the limit $a_0 \to \infty$, and with $n \ge d + 2$, the posterior covariance is $W_n^{-1}$, the inverse of a Wishart matrix, and the error is $\E \operatorname{tr} W_n^{-1} = d/(n - d - 1)$ \citep{muirhead1982}. Chapter~\ref{ch:specificity} builds its limits for a reliable description (Proposition~\ref{prop:reliable}) from these two facts.

A predictor that gives a density $g(y \mid x, D_n)$ rather than a number can instead be scored by its log loss, $-\log g(y \mid x, D_n)$. The expected log loss is smallest for the true conditional density of $y$ given the prompt \citep{cover2006}, which is the posterior predictive when tasks are drawn from the prior. Meta-trained networks are often trained with this loss (Section~\ref{sec:bg-bayes}), and so are the networks of Chapter~\ref{ch:networks}. The mean of the posterior predictive is the posterior mean, so a network that minimizes log loss and outputs a predictive mean can be scored by squared error. We use squared error throughout for two reasons. First, it needs only the learner's answer, a number, so it applies to a language model through the number that it writes (Chapter~\ref{ch:llm}). Second, the Bayes-optimal error~\eqref{eq:bg-gauss-regret} has the exact forms above, so the theory in Chapters~\ref{ch:specificity} and~\ref{ch:reliability} can work with exact errors.

Squared error also splits exactly over the components of a mixture. Suppose that a latent indicator $Z$ selects the component that generates the task. By the law of total variance, the posterior variance is the sum of two parts:
\begin{equation}\label{eq:bg-total-variance}
\begin{split}
    \operatorname{Var}(w^\top x \mid x, D_n) ={}& \E\big[\operatorname{Var}(w^\top x \mid x, D_n, Z) \mid x, D_n\big] \\
    &+ \operatorname{Var}\big(\E[w^\top x \mid x, D_n, Z] \mid x, D_n\big).
\end{split}
\end{equation}
The first part is the average variance within the components. The second is the variance of the prediction across the components. Chapter~\ref{ch:reliability} applies this decomposition to a description that may be irrelevant. There, the first part gives the risk of the Bayes-optimal predictor that also observes $Z$, and the second is the excess Bayes risk due to unobserved relevance.

\section{The Effective Sample Size of a Prior}\label{sec:bg-ess}
Bayesian statistics measures the information in a prior by its effective sample size (ESS), the number of observations that would carry the same information. In the conjugate normal model, this number is exact. Suppose that the observations are $y_i \sim \mathcal{N}(w, \sigma^2)$ and the prior is $w \sim \mathcal{N}(m, \sigma^2 / c)$. After $n$ observations with mean $\bar y$, the posterior is
\[
    w \mid y_{1:n} \sim \mathcal{N}\Big(\frac{c\,m + n \bar y}{c + n},\; \frac{\sigma^2}{c + n}\Big).
\]
This is the posterior that $c + n$ observations would give under a flat prior if $c$ of them averaged $m$. The prior is therefore worth $c$ observations, which is its precision measured in units of the precision of one observation. Chapter~\ref{ch:llm} uses this model for its prompt-based Bayesian reference, ignoring negligible boundary effects. The specificity $c$ of a description, defined in Chapter~\ref{ch:setup}, is the same conjugate-prior ESS, counted in each direction in units of direct observations of the task. Chapters~\ref{ch:specificity} and~\ref{ch:reliability} show how the ESS of a description, counted in demonstrations, differs from it.

\paragraph{Prior effective sample size.}
Prior ESS expresses the influence of a prior in observation units, with different definitions matching different statistical quantities \citep{morita2008,neuenschwander2020}. Our comparison is closest to measures based on expected loss \citep{reimherr2021,wiesenfarth2020}. \citet{lescao2021} measure the average horizontal separation of learning curves in prompt-based fine-tuning. \citet{selby2025} compare language-model-elicited priors with data-based predictors by the number of observations needed to match predictive loss. We adapt performance matching to task descriptions at a specified demonstration count, using each learner's own baseline curve, and analyze how precision, dimension, and relevance probability determine predictive ESS.

These definitions extend the ESS beyond the conjugate case in different ways. \citet{morita2008} extend it to any parametric prior. Their ESS is the number of hypothetical observations that, starting from a very vague prior, would give a posterior carrying about the same information as the prior, where information is measured by the curvature of the log density. \citet{neuenschwander2020} require an ESS to be predictively consistent: on average, the ESS of the posterior after $N$ observations should equal the prior's ESS plus $N$. Other definitions let the ESS depend on the current sample, through its size or through the observed data themselves. \citet{reimherr2021} compare a prior with a baseline prior. Their ESS is the number of observations that the baseline would need to perform as well at a given sample size, and it can be negative when the prior and the likelihood conflict. Similarly, the effective current sample size of \citet{wiesenfarth2020} counts the observations that must be added to, or removed from, the current sample for an analysis without the prior to match the analysis with it. Finally, \citet{reznik2026} measures information by predictive log loss to choose the discount of a power prior. A power prior raises the likelihood of historical data, which may not match the current data, to a power between $0$ and $1$.

Our predictive ESS (Definition~\ref{def:ess}) is the number of additional demonstrations that a learner without the description needs to match the learner's error with it. Like the measure of \citet{reznik2026}, it measures information by predictive loss, here squared error.

Our ESS also differs from that of \citet{morita2008} in what it matches. Their ESS matches information. In the conjugate normal model above, a prior of precision $c$ adds $c$ to the posterior precision whatever the number of observations, so it is worth $c$ observations at every sample size. In the same way, a description that raises the prior precision from $a_0^{-1}$ to $c$ in every direction is worth $c - a_0^{-1} = (1-r)\,c$ by their measure, whatever the number of demonstrations in hand. Our ESS matches errors instead. In $d$ dimensions, $n$ demonstrations leave the Bayesian predictor an error of $d/(n - d - 1)$ under a flat prior (Section~\ref{sec:bg-regret}), not the $d/n$ that their expected information would give. Matching errors therefore adds a term that grows with the dimension, the dimension term of Chapter~\ref{ch:specificity}. This term fades as demonstrations accumulate, and with many demonstrations in hand the $\ess_n$ of a reliable description returns to their value (Corollary~\ref{cor:ess}(ii)).

\section{Robust Mixture Priors}\label{sec:bg-mixture}
A prior built from historical data can conflict with the current data. \citet{evans2006} check for such prior--data conflict by asking whether the observed data are surprising under the prior predictive distribution. Clinical trials that borrow historical control data protect against conflict with a robust mixture prior \citep{schmidli2014}. The informative prior $f_I$, derived from the historical data, is mixed with a vague prior $f_V$:
\[
    f(w) = p\, f_I(w) + (1 - p)\, f_V(w).
\]
The posterior is then a mixture of the two component posteriors, and the weight on the informative component is updated in proportion to $p$ times its marginal likelihood. When the data agree with the informative component, its weight grows, and the trial borrows the historical information. When they conflict, the weight moves to the vague component, and the analysis falls back on the current data.

A description is a prior of the same kind. It may describe the task at hand, but it may also have been written for another task, or it may simply be wrong. Chapter~\ref{ch:setup} therefore models a description as a robust mixture prior whose weight is its reliability $p$, the probability that the description is relevant. Chapter~\ref{ch:reliability} calls the updated weight the posterior relevance probability $\pi_n$ and shows how reliability affects the description's ESS.

\section{Descriptions in Theories of ICL}\label{sec:bg-theory}
Additional task information has been studied through input-distribution descriptions \citep{huangge2025}, candidate hypotheses \citep{lin2025}, task-group mixtures \citep{lin2024}, changing instruction influence \citep{tong2026}, and related datasets supplied as context \citep{zhu2026}. We use explicit Bayesian predictors to quantify how such information changes the demonstrations needed for a given prediction error.

In some of these theories, the description is a set of tokens that indicate a hypothesis class \citep{lin2025} or carry the mean of the inputs \citep{huangge2025}. These theories generally do not measure the value of a description in demonstrations, and they do not model the possibility that it is irrelevant. For example, \citet{huangge2025} show that a learner can use the description token to predict better, but they do not say how many demonstrations it is worth or what it is worth when it may be wrong. The Bayesian theory of \citet{xie2022} models the prompt as demonstrations alone. We keep its view of the prompt as evidence about a latent task and add a description, which induces a distribution over tasks before the demonstrations are observed. Our networks in Chapter~\ref{ch:networks} use the prefix embedding of \citet{huangge2025}, in which the description is a single token placed before the demonstrations. The work closest to our setting is that of \citet{zhu2026}, who show that a Transformer given a prefix of related datasets performs amortized hierarchical Bayesian prediction. We measure the value of such conditioning information in demonstrations, and we allow it to be irrelevant.

\section{Descriptions and Priors in Pretrained Language Models}\label{sec:bg-llm}

Task summaries distilled from demonstrations can replace the examples in context \citep{honda2025}. Other studies report weak sensitivity to instruction relevance \citep{webson2022} or limited benefits from repository instruction files \citep{gloaguen2026}. \citet{gupta2025} compare explicitly stated coin biases with observations and evaluate the predictions against Bayesian references. Concurrent work studies how reasoning and code execution affect language models' approximation of Bayesian posteriors \citep{filbry2026}. Our numerical experiments vary stated precision and reliability separately, with controls that hold the description center and demonstrations fixed. Code execution is disabled in the main runs. Section~\ref{app:related} discusses further connections.

\section{Additional Related Work}\label{app:related}

\paragraph{Task descriptions and demonstrations.}
\citet{weller2020} formalize learning from task descriptions by distinguishing the latent task, its natural-language description, and individual input--output examples. We use \emph{task description} for information about the task supplied in the prompt, including a numeric description token in the meta-trained setting and probabilistic statements or policy constraints in the language-model experiments. A \emph{demonstration} is an observed input--output example; in scalar mean estimation, it is an observation at the fixed input $x=1$. Summaries distilled from demonstrations can serve as descriptions \citep{honda2025}; our controlled experiments instead specify the joint distribution of task, description, and demonstrations.

\paragraph{Prior information and uncertain relevance.}
The two-component construction is related to robust mixture priors, which combine informative and baseline distributions and update their component weights after observing data \citep{schmidli2014}. Our mixture is specified directly, rather than derived from a historical-data meta-analysis. Reliability $p$ is the prior probability of the description-conditioned component; $\pi_n$ is its posterior probability. Averaging component predictions using $\pi_n$ is Bayesian model averaging \citep{hoeting1999}. The experimental comparison that retains $p$ instead averages with prior weights. Fixed prior weights in a robust prior still yield updated posterior weights.

Our relevance terminology also has a conceptual precedent in information fusion, where relevance is distinguished from truthfulness \citep{pichon2012}. That work uses belief functions; our probability model is different. In our model, an irrelevant description is independent of the task and may happen to agree with it. Thus irrelevance is a generating condition, whereas prior--data conflict concerns disagreement between prior information and observed data. Related approaches discount external information according to predictive loss \citep{reznik2026} or accommodate conflict with heavy-tailed distributions \citep{ohagan2012}.

\paragraph{Specificity and description weights.}
Specificity takes a concrete form in each setting. In the isotropic Gaussian model, it is represented by prior precision $c$; a partial description also specifies the number $k$ of constrained directions. In the semantic experiment, broad and specific notes leave different numbers of compatible policies. These constructions do not define a common numerical scale across tasks. In Gaussian mean estimation, the reference combines the prior mean and sample mean with weights proportional to their precisions \citep{neuenschwander2020}. Fitted and per-response description weights summarize model predictions. Their normalized values can be compared with posterior component probabilities, but need not themselves be probabilities.

\paragraph{Effective sample size and predictive comparisons.}
Prior ESS has several definitions, including matching local information \citep{morita2008}, matching expected posterior loss \citep{reimherr2021}, and quantifying sample savings at a target estimation accuracy \citep{wiesenfarth2020}. Our definition follows the risk-matching approach but evaluates population prediction error and allows a non-Bayesian learner. Its dependence on the demonstration count is part of the definition. The word \emph{predictive} specifies what is matched; it does not assert the additive predictive-consistency property studied by \citet{neuenschwander2020}.

Two predictive comparisons are especially close. \citet{lescao2021} average horizontal data advantages over an accuracy range in supervised fine-tuning; we evaluate a pointwise difference in demonstration counts. \citet{selby2025} compare priors elicited from language models with data-based predictors by the observations needed to match predictive loss; we study descriptions supplied to a learner and how their benefit changes after demonstrations. Beyond ICL, an instruction can replace up to $100$ execution examples in the fine-tuning settings studied by \citet{cook2026}. \citet{wenliang2026} compare optimized prompts with demonstrations in terms of their length and effect. They treat a prompt as conditioning for a Bayesian sequence predictor and find that an optimized prompt can be shorter than demonstrations with the same effect.

\paragraph{Demonstrations, text, and probabilistic prediction.}
Other work examines semantic priors and flipped labels \citep{wei2023}, weak dependence on demonstration labels \citep{min2022}, Bayesian scaling laws \citep{arora2025}, and Bayesian reference curves for in-context function learning \citep{akata2026}. Text supplied alongside numbers can change numerical predictions \citep{requeima2024}, and models estimate a density from numbers in context like a kernel density estimator with an adaptive bandwidth \citep{liu2025}. Related work studies how language-model outputs compare with Bayesian inference and how training can improve probabilistic reasoning \citep{qiu2026}; \citet{he2025} examine Bayesian consistency during iterative reasoning on open-ended questions. Our numerical experiments state uncertainty explicitly, allowing comparison with a specified Bayesian reference as precision and relevance probability vary. The semantic task extends this comparison to ordinary-language constraints on a finite family of classification policies.

In the reference comparisons, the learning curves of pretrained models approach a Bayesian reference on synthetic tasks with a known optimum \citep{akata2026}, and \citet{arora2025} model in-context learning curves with a Bayesian scaling law that includes a task prior. Where \citet{min2022} and \citet{wei2023} vary the labels of the demonstrations, we ask the converse question: does a model weigh a description according to what the prompt says about it?
